% =================================================
% TikZ-рисунки к уроку 1
% =================================================

\newcommand{\digitcardsfigure}{%
  \begin{center}
    \begin{tikzpicture}[x=1cm,y=1cm]
      \foreach \d [count=\i from 0] in {0,1,2,3,4,5,6,7,8,9}{
        \node[
          draw=Primary,
          fill=PrimaryLight,
          rounded corners=1.4mm,
          minimum width=9mm,
          minimum height=10mm,
          font=\bfseries\sffamily\large,
          text=Primary
        ] at (\i*1.05,0) {\d};
      }
      \node[font=\small\sffamily,text=GrayText,anchor=north west]
        at (0,-0.7) {Из десяти цифр можно составить бесконечно много записей чисел.};
    \end{tikzpicture}
  \end{center}
}

\newcommand{\placevaluefigure}{%
  \begin{center}
    \begin{tikzpicture}[x=1.55cm,y=1cm]
      \foreach \x/\head/\digit in {
        0/{десятки\\тысяч}/5,
        1/{единицы\\тысяч}/8,
        2/{сотни}/3,
        3/{десятки}/4,
        4/{единицы}/2}{
        \node[
          draw=GrayLine,
          fill=TableHeader,
          minimum width=1.47cm,
          minimum height=1.35cm,
          align=center,
          font=\footnotesize\bfseries\sffamily,
          text=Primary
        ] at (\x,0.65) {\head};
        \node[
          draw=GrayLine,
          fill=white,
          minimum width=1.47cm,
          minimum height=1.15cm,
          font=\bfseries\sffamily\Large
        ] at (\x,-0.6) {\digit};
      }
      \node[font=\small\sffamily,text=GrayText,anchor=north]
        at (2,-1.35) {Число \(58\,342\): каждая цифра занимает свой разряд.};
    \end{tikzpicture}
  \end{center}
}

\newcommand{\placevalueworksheetfigure}{%
  \begin{center}
    \begin{tikzpicture}[x=1.55cm,y=1cm]
      \foreach \x/\head/\digit in {
        0/{десятки\\тысяч}/7,
        1/{единицы\\тысяч}/2,
        2/{сотни}/6,
        3/{десятки}/4,
        4/{единицы}/9}{
        \node[
          draw=GrayLine,
          fill=TableHeader,
          minimum width=1.47cm,
          minimum height=1.35cm,
          align=center,
          font=\footnotesize\bfseries\sffamily,
          text=Primary
        ] at (\x,0.65) {\head};
        \node[
          draw=GrayLine,
          fill=white,
          minimum width=1.47cm,
          minimum height=1.15cm,
          font=\bfseries\sffamily\Large
        ] at (\x,-0.6) {\digit};
      }
    \end{tikzpicture}
  \end{center}
}

\newcommand{\zeroplacefigure}{%
  \begin{center}
    \begin{tikzpicture}[node distance=8mm]
      \node[
        draw=Primary,fill=PrimaryLight,rounded corners=1.5mm,
        minimum width=4.6cm,minimum height=13mm,
        font=\bfseries\sffamily\Large,text=Primary
      ] (n) {40\,508};
      \node[
        below=5mm of n,
        font=\small\sffamily,
        text=GrayText,
        align=center
      ] {0 тысяч и 0 десятков: нули сохраняют места разрядов};
      \draw[-{Stealth[length=2.5mm]},draw=Accent,line width=0.8pt]
        ($(n.center)+(-0.48,0.05)$) -- ++(0,-0.62);
      \draw[-{Stealth[length=2.5mm]},draw=Accent,line width=0.8pt]
        ($(n.center)+(0.48,0.05)$) -- ++(0,-0.62);
    \end{tikzpicture}
  \end{center}
}

\newcommand{\expandedformfigure}{%
  \begin{center}
    \begin{tikzpicture}[node distance=7mm]
      \node[
        draw=Primary,fill=PrimaryLight,rounded corners=1.5mm,
        inner xsep=7mm,inner ysep=3.5mm,
        font=\bfseries\sffamily\Large,text=Primary
      ] (number) {47\,305};
      \node[
        below=8mm of number,
        draw=Secondary,fill=SecondaryLight,rounded corners=1.5mm,
        inner xsep=5mm,inner ysep=3.5mm,
        font=\sffamily
      ] (sum) {\(40\,000+7\,000+300+5\)};
      \draw[-{Stealth[length=3mm]},draw=Accent,line width=1pt]
        (number) -- node[right,font=\small\sffamily,text=GrayText]{разложили по разрядам} (sum);
    \end{tikzpicture}
  \end{center}
}
